%% Separately authored additions after exact Modal, Temporal, and Epistemic restoration.
% Necessity: the exact inherited chapter does not contain the LTL system card,
% branching-time, fixed-point, automata, tense, doxastic, dynamic, or comparison
% material below, so those separately attributed additions would otherwise
% vanish from the reader route.
% Definition: this editorial derivative descends from four historical ranges
% and includes current semantic, system-card, source-role, and apparatus
% repairs; it is not a literal reproduction of those ranges.
% Standing: it is editorial reader material, not inherited source or a silent
% proof-status promotion.
% Nonredundancy: the exact wrappers provide the complete source chapter, inline
% proof sketches, examples, and exercises; these derivatives retain only
% distinct later material while both historical bodies and annexes stay pinned.
% Use: parts/p08.tex inputs this file exactly once after the modal supplement;
% route, slice, label, consumer, and publication tests fail closed on drift.
\noindent\textit{Reader note.} The following temporal and epistemic cards,
comparisons, and later semantics are separately attributed additions to
the exact inherited chapter above; they neither replace nor alter its claims.

% BEGIN HISTORICAL DERIVATIVE 49-68
\begin{systementry}{entry:p08-ltl}{Linear temporal logic LTL}
  \index{linear temporal logic}
  \entrylanguage{The infinite-word language of
    Section~\ref{sec:modal-temporal-epistemic-02}, with next and strong until primitive.}
  \entrysemantics{Truth at a position of a discrete infinite word; validity
    quantifies over every word and every position.}
  \entrycalculus{sheet:hilbert-ltl}
  \entryadequacy{sketched=annex:p8.ltl.source}{Sound and weakly complete
    for validity.  The exact inherited source gives the theorem and inline proof sketch in
    Section~\ref{sec:modal-temporal-epistemic-02}; no arbitrary-premise
    strong-completeness claim is made.}
  \entrydecidability{cited=SistlaClarke1985}{Satisfiability and combined
    finite-state model checking are PSPACE-complete.}
  \entrycharacteristic{Expresses linear safety, liveness, and recurrence;
    finite traces and branching time use different semantics.}
  \entrytranslations{Every formula has an equivalent nondeterministic
    B\"uchi automaton; Figure~\ref{fig:p08-buchi-gfp} gives one exact worked
    instance, not the general translation.}
  \entrysources{GabbayPnueliShelahStavi1980,Pnueli1977,SistlaClarke1985,VardiWolper1986}
\end{systementry}
% END HISTORICAL DERIVATIVE 49-68

% BEGIN HISTORICAL DERIVATIVE 89-508
\section{Branching time and fixed points}

\begin{hbdefinition}{def:p08-transition-structure}{Total transition structure and path}
A \emph{total transition structure} is
\(\mathcal K=(S,R,V)\), where \(S\ne\varnothing\),
\(R\subseteq S\times S\), every state has an \(R\)-successor, and
\(V:\mathsf{Prop}\to\mathcal P(S)\).  A \emph{path from \(s\)} is a map
\(\pi:\mathbb N\to S\) with \(\pi(0)=s\) and
\(\pi(n)R\pi(n+1)\) for every \(n\).  Totality is necessary here so every
path is infinite.  A state is a semantic evaluation point; a path is a
sequence of such points.  Their different standings are used immediately to
separate CTL state formulas from CTL* path formulas.
\end{hbdefinition}

\begin{hbdefinition}{def:p08-ctl-semantics}{CTL state formulas}
CTL formulas form the least Boolean-closed set of state formulas containing
the atoms and closed under
\(\mathsf{EX}A\), \(\mathsf{AX}A\),
\(\mathsf E(A\mathbin{\mathsf U}B)\), and
\(\mathsf A(A\mathbin{\mathsf U}B)\).  The letters \(\mathsf E\) and
\(\mathsf A\) are path quantifiers, not epistemic operators: at a state \(s\),
\(\mathsf E\) selects some path from \(s\) and \(\mathsf A\) quantifies over
all paths from \(s\).  Thus
\[
\begin{aligned}
 \mathcal K,s\entails\mathsf{EX}A
 &\text{ iff }
   \exists t\,(sRt\land\mathcal K,t\entails A),\\
 \mathcal K,s\entails\mathsf{AX}A
 &\text{ iff }
   \forall t\,(sRt\to\mathcal K,t\entails A),
\end{aligned}
\]
 and the strong-until clauses are
 \[
 \begin{aligned}
  \mathcal K,s\entails\mathsf E(A\mathbin{\mathsf U}B)
  &\text{ iff }
    \exists\pi\,\exists m\in\mathbb N\,
    \left(
      \begin{gathered}
        \pi\text{ is a path from }s\\
        {}\land\ \mathcal K,\pi(m)\entails B\\
        {}\land\ \forall k<m\,\mathcal K,\pi(k)\entails A
      \end{gathered}
    \right),\\
  \mathcal K,s\entails\mathsf A(A\mathbin{\mathsf U}B)
  &\text{ iff }
    \forall\pi\,
    \left(
      \begin{gathered}
        \pi\text{ is a path from }s\ \to\ \exists m\in\mathbb N\,\\
        \bigl(\mathcal K,\pi(m)\entails B\ \land\
          \forall k<m\,\mathcal K,\pi(k)\entails A\bigr)
      \end{gathered}
    \right).
 \end{aligned}
 \]
 The eventuality abbreviations used below are
 \[
  \mathsf{EF}A:=\mathsf E(\top\mathbin{\mathsf U}A),
  \qquad
  \mathsf{AF}A:=\mathsf A(\top\mathbin{\mathsf U}A),
 \]
 where \(\top\) is the Boolean truth constant.  These are abbreviations
 for the displayed CTL state constructors, whose until witnesses allow
 \(m=0\).
 This paired syntax and its state-level clauses distinguish CTL from both LTL
 and CTL* and are used by the card and comparison table below.
\end{hbdefinition}

\begin{systementry}{entry:p08-ctl}{Computation tree logic CTL}
  \index{computation tree logic!CTL}
  \entrylanguage{Boolean state formulas with path quantifiers immediately
    paired with next or strong until, as in
    Definition~\ref{def:p08-ctl-semantics}.}
  \entrysemantics{States of total transition structures; path quantifiers range
    over infinite paths beginning at the current state.}
  \entrycalculus{sheet:hilbert-ctl}
  \entryadequacy{cited=EmersonHalpern1985}{The registered theoremhood calculus is
    sound and weakly complete for CTL validity on total transition structures.}
  \entrydecidability{cited=EmersonHalpern1985}{Satisfiability is
    EXPTIME-complete; combined finite-state model checking is P-complete and
    admits the standard linear-time-in-structure-times-formula algorithm
    \cite{ClarkeEmersonSistla1986}.}
  \entrycharacteristic{Can distinguish existential from universal branching,
    as Figure~\ref{fig:p08-ctl-tree} illustrates, but cannot retain an arbitrary
    LTL path formula beneath one path quantifier.}
  \entrytranslations{Embeds compositionally into CTL*.  Its fixed-point
    translation is recorded below only after the target language has been
    defined.}
  \entrysources{ClarkeEmerson1981,EmersonHalpern1985,ClarkeEmersonSistla1986}
\end{systementry}

\begin{hbdefinition}{def:p08-ctl-star-semantics}{The state/path distinction in CTL*}
CTL* uses two mutually generated syntactic sets.  State formulas contain atoms,
Boolean combinations, and \(\mathsf A\alpha\) when \(\alpha\) is a path
formula.  Here \(\alpha,\beta\) range over path formulas.  Path formulas
contain every state formula and are closed under Boolean connectives,
\(\mathsf X\alpha\), and \(\alpha\mathbin{\mathsf U}\beta\).  Define
\(\mathsf E\alpha\defeq\neg\mathsf A\neg\alpha\).  The bridge from state
truth to path truth and the path-quantifier clause are
\[
 \begin{aligned}
  \mathcal K,\pi,n\entails A
  &\text{ iff } \mathcal K,\pi(n)\entails A,
  &&\text{for every state formula (A)},\\
  \mathcal K,s\entails\mathsf A\alpha
  &\text{ iff }
    \forall\pi\,(\pi\text{ is a path from }s
      \to\ \mathcal K,\pi,0\entails\alpha).
 \end{aligned}
\]
On path formulas,
\[
 \begin{aligned}
  \mathcal K,\pi,n\entails\mathsf X\alpha
  &\text{ iff } \mathcal K,\pi,n+1\entails\alpha,\\
  \mathcal K,\pi,n\entails\alpha\mathbin{\mathsf U}\beta
  &\text{ iff }
    \exists m\ge n\,
      \bigl(\mathcal K,\pi,m\entails\beta\ \land\
      \forall k\,(n\le k<m\to\mathcal K,\pi,k\entails\alpha)\bigr).
 \end{aligned}
\]
Boolean clauses are classical at the appropriate state or path coordinate.
The separate syntactic standings are
necessary because CTL* permits an arbitrary temporal path formula below a path
quantifier; this is precisely what CTL forbids and what the card uses.
\end{hbdefinition}

\begin{systementry}{entry:p08-ctl-star}{Full computation tree logic CTL*}
  \index{computation tree logic!CTL*}
  \entrylanguage{The mutually generated state/path language of
    Definition~\ref{def:p08-ctl-star-semantics}.}
  \entrysemantics{State formulas are evaluated at states and path formulas on
    infinite paths of total transition structures.}
  \entrycalculus{sheet:hilbert-ctl-star}
  \entryadequacy{cited=Reynolds2001}{Reynolds's fullpath system, including its
    limit-closure schema and exact auxiliary-atoms side condition, is sound and
    weakly complete for standard CTL* validity.  Noncompactness precludes the
    corresponding arbitrary-premise strong-completeness claim.}
  \entrydecidability{cited=EmersonJutla1999}{Satisfiability is
    2EXPTIME-complete; combined finite-state model checking is PSPACE-complete
    \cite{EmersonLei1987}.}
  \entrycharacteristic{Strictly contains CTL and state-lifted LTL.  Its
    comparison with the later fixed-point language is stated only after that
    target language is defined.}
  \entrytranslations{The later fixed-point card records the exact target of
    the standard translation, generally with a nontrivial size increase.}
  \entrysources{EmersonHalpern1986,Reynolds2001,EmersonLei1987,EmersonJutla1999}
\end{systementry}

\begin{hbdefinition}{def:p08-modal-mu-semantics}{Kozen's \(L_{\mu}\) syntax and semantics}
Use the finitary language introduced in \textcite{Kozen1983}.  Its primitive
nonlogical symbols are propositional constants \(p,q,\ldots\), including
\(\bottom\) and \(\topv\), fixed-point variables \(Z,Y,\ldots\), and program
constants \(a,b,\ldots\).  Formulas are the least set generated by
\[
 Z,\qquad p,\qquad A\lor B,\qquad \neg A,\qquad
 \langle a\rangle A,\qquad \mu Z.A,
\]
where, in \(\mu Z.A\), every free occurrence of \(Z\) in \(A\) lies in the
scope of an even number of negations.  This is Kozen's primitive grammar, not
a positive-normal-form or Hilbert-style modal replacement.  The
usual Boolean operations and the remaining modal and fixed-point operations
are defined on that grammar by
\[
 \begin{aligned}
 A\land B&\defeq\neg(\neg A\lor\neg B),
 & A\to B&\defeq\neg A\lor B,\\
 [a]A&\defeq\neg\langle a\rangle\neg A,
 & \nu Z.A&\defeq\neg\mu Z.\neg A[\neg Z/Z].
 \end{aligned}
\]
The substitutions are capture-avoiding.  Positive normal form is a derived
normal form for this language; it is not the grammar of the calculus.

A Kozen model is \(\mathcal M=(S,\mathcal I)\), where \(S\ne\varnothing\),
\(\mathcal I(p)\subseteq S\), \(\mathcal I(\bottom)=\varnothing\),
\(\mathcal I(\topv)=S\), and
\(\mathcal I(a)\subseteq S\times S\).  An environment \(\zeta\) maps
fixed-point variables to subsets of \(S\).  Define the metalinguistic
formula-extension map \(\operatorname{Ext}_{\zeta}(A)\subseteq S\) on the
primitive grammar by
\[
\begin{aligned}
 \operatorname{Ext}_{\zeta}(Z)&=\zeta(Z),\\
 \operatorname{Ext}_{\zeta}(p)&=\mathcal I(p),\\
 \operatorname{Ext}_{\zeta}(A\lor B)&=
   \operatorname{Ext}_{\zeta}(A)\cup\operatorname{Ext}_{\zeta}(B),\\
 \operatorname{Ext}_{\zeta}(\neg A)&=
   S\setminus\operatorname{Ext}_{\zeta}(A),\\
 \operatorname{Ext}_{\zeta}(\langle a\rangle A)&=
   \{s:\exists t\,((s,t)\in\mathcal I(a)\land
      t\in\operatorname{Ext}_{\zeta}(A))\},\\
 \operatorname{Ext}_{\zeta}(\mu Z.A)
 &=\bigcap\{X\subseteq S:
     \operatorname{Ext}_{\zeta[Z\mapsto X]}(A)\subseteq X\}.
\end{aligned}
\]
The name is needed to state fixed points as sets of states without confusing
formula extension with term denotation; its standing is exactly this
eliminable recursion.  The defined greatest fixed point consequently satisfies
\[
 \operatorname{Ext}_{\zeta}(\nu Z.A)
 =\bigcup\{X\subseteq S:
    X\subseteq\operatorname{Ext}_{\zeta[Z\mapsto X]}(A)\}.
\]
These are respectively least and greatest fixed-point extensions, not
unbounded minimization or lambda binding.  Positivity makes the induced
powerset operator monotone.  The distinction is used by Kozen's exact
equational sheet and by the expressiveness and model-checking fields below.
\end{hbdefinition}

\begin{systementry}{entry:p08-modal-mu}{Propositional modal \(\mu\)-calculus}
  \index{modal mu-calculus@modal \(\mu\)-calculus}
  \entrylanguage{Kozen's primitive \(L_{\mu}\) language of
    Definition~\ref{def:p08-modal-mu-semantics}.}
  \entrysemantics{State extensions in Kozen models with one relation per
    program constant; least fixed points are taken in \(\mathcal P(S)\), and
    greatest fixed points are the displayed duals.}
  \entrycalculus{sheet:hilbert-mu}
  \entryadequacy{cited=Walukiewicz2000}{Kozen's equational axiomatization is
    sound and weakly complete for validity.  The exact theorem, equation
    form, finite-premise consequence, and failure of compactness are in
    Section~\ref{sec:p09-kozen-full-completeness}.}
  \entrydecidability{cited=Kozen1983}{Satisfiability is EXPTIME-complete.
    Finite-state model checking is polynomial-time equivalent to solving parity
    games and has a quasipolynomial algorithm; a polynomial-time algorithm is
    not known \cite{CaludeJainKhoussainovLiStephan2017}.}
  \entrycharacteristic{Captures bisimulation-invariant modal fixed-point
    properties; its one-program fragment strictly extends CTL* on transition
    systems.}
  \entrytranslations{CTL, LTL under state lifting, and CTL* embed into the
    one-program fragment by using one program constant for the transition
    relation.}
  \entrysources{Kozen1983,Walukiewicz2000,EmersonJutla1999,CaludeJainKhoussainovLiStephan2017}
\end{systementry}

\begin{table}[!htbp]
\centering
\scriptsize
\begin{tabular}{@{}p{.08\textwidth}p{.20\textwidth}p{.23\textwidth}p{.18\textwidth}p{.21\textwidth}@{}}
\toprule
Logic & Semantic scope & Expressiveness & Satisfiability & Model checking \\
\midrule
LTL & Path formulas on infinite words & Embeds into CTL* after an outer path
quantifier; not itself a branching-state language & PSPACE-complete &
PSPACE-complete combined; linear in the structure for a fixed formula \\
CTL & State formulas with paired path/temporal operators & Strict fragment of
CTL* and modal \(\mu\) & EXPTIME-complete & P-complete; linear in structure
times formula size \\
CTL* & Separate state and path formulas & Strictly contains CTL and state-lifted
LTL; strict fragment of modal \(\mu\) & 2EXPTIME-complete & PSPACE-complete \\
Modal \(\mu\) & Bisimulation-invariant fixed-point state formulas & Strictly
extends CTL* on transition systems & EXPTIME-complete & Parity-game equivalent;
P-hard, in \(\mathsf{UP}\cap\mathsf{coUP}\), and quasipolynomial-time; P open \\
\bottomrule
\end{tabular}
\caption{Temporal-logics comparison.  Complexity is combined
complexity on finite, explicitly represented total transition structures.
The LTL row asks whether all initial paths satisfy the path formula; changing
the input convention changes the model-checking column
\cite{SistlaClarke1985,ClarkeEmersonSistla1986,EmersonLei1987,
EmersonJutla1999,CaludeJainKhoussainovLiStephan2017}.}
\label{tab:p08-temporal-comparison}
\end{table}

\section{Worked temporal diagrams}

\begin{figure}[!htbp]
\centering
\begin{tikzpicture}[>=Latex,font=\small,x=1.6cm,y=1cm]
  \foreach \x/\lab in {0/{0:\ p},1/{1:\ p},2/{2:\ p},3/{3:\ q},4/{4:\ p}}
    \node[draw,circle,minimum size=9mm] (n\x) at (\x,0) {\lab};
  \draw[->] (n0) -- (n1);
  \draw[->] (n1) -- (n2);
  \draw[->] (n2) -- (n3);
  \draw[->] (n3) -- (n4);
  \draw[->] (n4) -- ++(1.1,0) node[right]{\(p,p,p,\ldots\)};
  \node[draw=none,align=left,below=8mm of n2] {
    \(p\mathbin{\mathsf U}q\) is true at \(0,1,2,3\) and false at \(4\);\\
    at \(0\), \(\mathsf Xp\), \(\mathsf Fq\), and
    \(\mathsf G(p\lor q)\) are true, while \(\mathsf Gp\) is false};
\end{tikzpicture}
\caption{LTL evaluation on one infinite path.  The five
displayed truth values of \(p\mathbin{\mathsf U}q\) answer the inherited
position-by-position calculation.  The labelled tail is part of the structure;
the five displayed positions alone would not determine the two
\(\mathsf G\)-claims.}
\label{fig:p08-ltl-path}
\end{figure}

\begin{figure}[!htbp]
\centering
\begin{tikzpicture}[>=Latex,font=\small,node distance=1.4cm and 2.7cm]
  \node[draw,circle] (w) {\(w:\neg p\)};
  \node[draw,circle,below left=of w] (u) {\(u:p\)};
  \node[draw,circle,below right=of w] (v) {\(v:\neg p\)};
  \node[draw,circle,below=of u] (u1) {\(u_1:p\)};
  \node[draw,circle,below=of v] (v1) {\(v_1:\neg p\)};
  \node[draw=none,below=7mm of u1] (udots) {\(\vdots\ p\)};
  \node[draw=none,below=7mm of v1] (vdots) {\(\vdots\ \neg p\)};
  \draw[->] (w) -- (u);
  \draw[->] (w) -- (v);
  \draw[->] (u) -- (u1);
  \draw[->] (v) -- (v1);
  \draw[->] (u1) -- (udots);
  \draw[->] (v1) -- (vdots);
  \node[draw=none,align=left,right=18mm of v]
    {at \(w\): \(\mathsf{EF}p\) true, \(\mathsf{AF}p\) false;\\
     \(\mathsf{EX}p\) true, \(\mathsf{AX}p\) false};
\end{tikzpicture}
\caption{A total CTL computation tree contrasting existential and universal
path quantification.  The dots denote infinite continuing
branches, not terminal leaves.}
\label{fig:p08-ctl-tree}
\end{figure}

\begin{hbdefinition}{def:p08-buchi-automaton}{B\"uchi automaton and acceptance}
A \emph{nondeterministic B\"uchi word automaton} over an alphabet \(\Sigma\)
is a tuple
\(\mathcal A=(Q,\Sigma,\delta,q_0,F)\), where \(Q\) is finite,
\(q_0\in Q\), \(F\subseteq Q\), and
\(\delta:Q\times\Sigma\to\mathcal P(Q)\).  A run on
\(\sigma_0\sigma_1\ldots\) is a sequence \(r_0r_1\ldots\) with
\(r_0=q_0\) and \(r_{n+1}\in\delta(r_n,\sigma_n)\) for every \(n\).
Define the metalinguistic recurrence set
\[
 \operatorname{Inf}(r)\defeq
 \{q\in Q:\forall n\,\exists m\ge n\ (r_m=q)\}.
\]
The word is accepted exactly when there exists a run \(r\) with
\(\operatorname{Inf}(r)\cap F\ne\varnothing\).  The automaton is
\emph{deterministic} when every transition set is a singleton; its unique run
takes \(r_{n+1}\) to the unique member of
\(\delta(r_n,\sigma_n)\), as in the figure below.
The nondeterministic qualifier is necessary because the LTL translation on the
system card uses existential choice among runs; deterministic B\"uchi
automata are a proper specialization.  The recurrence set is needed to
state infinite recurrence without confusing it with the range of the run; it
differs from the latter by retaining only states visited infinitely often and
is used immediately in the acceptance test.  The test distinguishes B\"uchi
acceptance from a finite-word final state and is used by
Figure~\ref{fig:p08-buchi-gfp}
\cite{Buchi1962,VardiWolper1986}.
\end{hbdefinition}

\begin{figure}[!htbp]
\centering
\begin{tikzpicture}[>=Latex,font=\small,node distance=4cm]
  \node[draw,circle,minimum size=13mm] (qn) {\(q_{\neg p}\)};
  \node[draw,double,circle,minimum size=13mm,right=of qn] (qp) {\(q_p\)};
  \draw[->] ($(qn)+(-1.35,0)$) -- (qn) node[midway,above]{start};
  \draw[->,bend left=18] (qn) to node[above]{\(p\)} (qp);
  \draw[->,bend left=18] (qp) to node[below]{\(\neg p\)} (qn);
  \draw[->] (qn) to[out=225,in=135,looseness=6] node[left]{\(\neg p\)} (qn);
  \draw[->] (qp) to[out=45,in=-45,looseness=7] node[right]{\(p\)} (qp);
\end{tikzpicture}
\caption{The total deterministic B\"uchi automaton for
\(\mathsf G\mathsf Fp\).  Here
\(Q=\{q_{\neg p},q_p\}\), \(\Sigma=\mathcal P(\{p\})\), the initial state is
\(q_{\neg p}\), and \(F=\{q_p\}\).  From either state a letter containing
\(p\) leads to \(q_p\), and a letter omitting \(p\) leads to
\(q_{\neg p}\).  Hence the unique run visits the accepting state infinitely
often exactly when \(p\) occurs infinitely often, which is exactly the
infinite-word truth condition for \(\mathsf G\mathsf Fp\).}
\label{fig:p08-buchi-gfp}
\end{figure}

\section[Philosophical tense logic]{Philosophical tense logic and modal combinations}

\begin{hbdefinition}{def:p08-tense-until-since}{Kamp's strict U and S operators}
Let \((T,<,V)\) be a linearly ordered temporal structure.  The strict binary
future and past constructors have the clauses
\[
\begin{aligned}
 t\entails U\,A\,B
 &\text{ iff }
 \exists u>t\,[u\entails A\land
   \forall v\,(t<v<u\to v\entails B)],\\
 t\entails S\,A\,B
 &\text{ iff }
 \exists u<t\,[u\entails A\land
   \forall v\,(u<v<t\to v\entails B)].
\end{aligned}
\]
The prefix symbols \(U\) and \(S\) are Kamp's binary future and past formula
constructors, respectively.  Their first operand holds at the witness and
their second operand at every intervening point; they are not metalinguistic
English tense.  They
are needed because Prior's unary future/past operators cannot express all
one-variable first-order order properties.  Their strict witnesses and two
directions distinguish them from the non-strict discrete LTL until fixed in
Section~\ref{sec:modal-temporal-epistemic-02}; the following theorem uses them
\cite{Prior1967,Kamp1968}.
\end{hbdefinition}

\begin{hbtheorem}{thm:p08-kamp}{cited=Kamp1968}{Kamp expressive completeness with its order hypothesis}
Kamp's compositional map sends every formula using negation, conjunction, and
his strict \(U,S\) to an equivalent one-free-variable formula in the displayed
first-order language with \(<\) and unary predicates; its truth induction is
valid over every binary structure.  Conversely, over Kamp's class
\(\mathsf{Com}\) of strict linear orders, every uniformly
first-order-definable tense is expressible by such a temporal formula.
\(\mathsf{Com}\) requires greatest lower and least upper bounds for every
nonempty bounded set; it permits discrete orders and endpoints and does not
mean deductive completeness.  The reverse translation is primitive recursive,
but the thesis gives no complexity or output-size bound.  It is not asserted
over all linear orders; Kamp's rational-order Theorem~IV.5 supplies a
counterexample to this four-operation basis.  The full source presentation is
Section~\ref{sec:p09-presentation-n-c-012}.
\end{hbtheorem}

Prior's unary \(\mathsf F,\mathsf G\) future pair and
\(\mathsf P,\mathsf H\) past pair can be interpreted through a future order
and its converse.  Different Prior systems add assumptions such as
transitivity, linearity, density, or endpoints; this S-level survey does not
select a further calculus or transfer an adequacy theorem between those frame
classes \cite{Prior1967,Burgess1984}.

\begin{hbdefinition}{def:p08-modal-fusion}{Fusion of normal modal logics}
For normal modal logics \(L_1,L_2\) in disjoint modal signatures, their
\emph{fusion} is the least normal multimodal logic containing the theorem sets
of both components and no interaction axioms.  Its standing is a construction
on logics, not a connective.  The name is necessary to distinguish mere
 side-by-side combination from products, commutation principles, and tense
 systems whose past and future relations are converses; that nonredundancy is
 used in the following comparison.  In particular, placing independent future
 and past logics side by side gives their fusion, whereas requiring their frame
 relations to be converses adds an interaction condition and no longer gives a
 pure fusion.
\end{hbdefinition}

\begin{hbdefinition}{def:p08-ockhamist-peircean}{Ockhamist and Peircean branching-time evaluation}
On a branching order, a \emph{history} is a maximal linearly ordered set of
moments.  Ockhamist future-tense truth is evaluated at a moment--history pair,
so it can depend on the current history through that moment.  Peircean
future-tense truth is evaluated at a moment and quantifies over branching
futures.  These are semantic evaluation policies, not competing names for CTL
path quantifiers.  Their different evaluation points prevent ambiguity, and
the distinction is used to delimit the computational cards above from this
philosophical survey \cite{Burgess1984}.
\end{hbdefinition}

% END HISTORICAL DERIVATIVE 89-508

% BEGIN HISTORICAL DERIVATIVE 526-564
\begin{systementry}{entry:p08-s5n}{Multi-agent epistemic logic \(\mathsf{S5}_{I}\)}
  \index{epistemic logic!multi-agent S5}
  \entryaliases{\(\mathsf{S5}_n\) when \(I=\{1,\ldots,n\}\)}
  \entrylanguage{Classical propositional formulas with one knowledge operator
    \(K_i\) for each \(i\in I\).}
  \entrysemantics{Epistemic structures whose \(\sim_i\) relations are
    equivalence relations, with no imposed interaction between distinct agents.}
  \entrycalculus{sheet:hilbert-s5n}
  \entryadequacy{sketched=annex:p8.s5n.source}{Sound and strongly complete
    for local consequence.  The exact inherited source gives the theorem and inline proof sketch in
    Section~\ref{sec:modal-temporal-epistemic-03}.}
  \entrydecidability{cited=FaginHalpernMosesVardi1995}{The finite-agent
    satisfiability problem is decidable; no sharper bound is required by this
    source-backed card.}
  \entrycharacteristic{Factivity and positive and negative introspection hold
    componentwise; there are no cross-agent interaction axioms.}
  \entrytranslations{Group operators are defined below from the indexed
    relations; they are not silently built into this base card.}
  \entrysources{Hintikka1962,FaginHalpernMosesVardi1995}
\end{systementry}

\begin{systementry}{entry:p08-kd45n}{Multi-agent doxastic logic \(\mathsf{KD45}_{I}\)}
  \index{doxastic logic!multi-agent KD45}
  \entryaliases{\(\mathsf{KD45}_n\) when \(I=\{1,\ldots,n\}\)}
  \entrylanguage{Classical propositional formulas with one belief operator
    \(B_i\) for each \(i\in I\).}
  \entrysemantics{One serial, transitive, and Euclidean accessibility relation
    per agent, with no cross-agent interaction axioms.}
  \entrycalculus{sheet:hilbert-kd45n}
  \entryadequacy{cited=FaginHalpernMosesVardi1995}{Sound and strongly complete
    for local consequence on the stated componentwise frame class.}
  \entrydecidability{cited=FaginHalpernMosesVardi1995}{The finite-agent
    satisfiability problem is decidable.}
  \entrycharacteristic{Consistency and positive and negative introspection
    hold, but factivity need not: this is a belief rather than knowledge card.}
  \entrytranslations{Replacing every equivalence relation by a merely serial,
    transitive, Euclidean one changes both terminology and frame class.}
  \entrysources{Hintikka1962,FaginHalpernMosesVardi1995}
\end{systementry}
% END HISTORICAL DERIVATIVE 526-564

% BEGIN HISTORICAL DERIVATIVE 610-774
\begin{hbdefinition}{def:p08-distributed-knowledge}{Distributed knowledge for a finite group}
Let \(G\subseteq I\) be nonempty and finite.  In the multi-agent epistemic
model with agent relations \(\sim_i\), define
\[
 R_G^{D}\defeq\bigcap_{i\in G}\sim_i,
 \qquad
 \mathcal M,w\models D_GA
 \quad\Longleftrightarrow\quad
 \forall u\,(wR_G^{D}u\Rightarrow\mathcal M,u\models A).
\]
The name is necessary because pooled information can determine \(A\) even
when no individual agent knows it.  It names a group modality, not a new frame
relation for an individual agent.  It differs from
\(E_GA\defeq\bigwedge_{i\in G}K_iA\), which uses the union of one-step agent
alternatives, and from \(C_GA\), which uses the reflexive--transitive closure
of that union.  The distinction is used by the card and comparison table below
\cite{FaginHalpernMosesVardi1995}.
\index{distributed knowledge}\index{group knowledge}
\end{hbdefinition}

\begin{systementry}{entry:p08-common-knowledge}{Multi-agent \(\mathsf{S5}\) with common knowledge}
  \index{common knowledge!system entry}
  \entryaliases{\(\mathsf{S5}_{n}^{C}\) when \(I=\{1,\ldots,n\}\)}
  \entrylanguage{The \(\mathsf{S5}_{I}\) language extended by \(C_G\) for
    nonempty finite groups, with \(E_G\) defined as in
    Section~\ref{sec:modal-temporal-epistemic-03}.}
  \entrysemantics{\(C_G\) quantifies over the reflexive-transitive closure of
    the union of the agents' equivalence relations.}
  \entrycalculus{sheet:hilbert-common-knowledge}
  \entryadequacy{cited=FaginHalpernMosesVardi1995}{Sound and weakly complete
    for validity.  The exact fixed-point and induction surfaces are on the
    sheet; arbitrary-premise strong completeness fails by noncompactness.}
  \entrydecidability{cited=FaginHalpernMosesVardi1995}{The standard
    finite-agent common-knowledge language is decidable.}
  \entrycharacteristic{Common knowledge is not any fixed finite iteration of
    everyone-knows; its semantic closure contains paths of every finite length.}
  \entrytranslations{It is a greatest-fixed-point modality, while
    \(D_G\) is Definition~\ref{def:p08-distributed-knowledge} and remains a
    separate intersection-based operator.}
  \entrysources{Aumann1976,FaginHalpernMosesVardi1995}
\end{systementry}

\begin{hbexample}{def:p08-everyone-not-common}{Everyone knows without common knowledge}
Let \(G=\{i,j\}\), let \(W=\{w,u,v\}\), and make the \(i\)-equivalence
classes \(\{w,u\},\{v\}\) and the \(j\)-equivalence classes
\(\{w\},\{u,v\}\).  Let \(p\) hold exactly at \(w,u\).  At \(w\), every
\(i\)-accessible world and every \(j\)-accessible world satisfies \(p\), so
\(E_Gp\) holds.  But \(w\sim_i u\sim_j v\), and \(p\) is false at \(v\);
hence \(C_Gp\) fails at \(w\).  This realizes the source's third epistemic
exercise and exhibits why finite conjunction is weaker than reachability.
\end{hbexample}

\begin{hbtheorem}{thm:p08-common-knowledge-noncompact}{cited=FaginHalpernMosesVardi1995}{Common knowledge and noncompactness}
Let \(G\) contain at least two agents, put \(E_G^0p=p\), and put
\(E_G^{n+1}p=E_GE_G^np\).  Then
\[
 \{E_G^np:n\in\mathbb N\}\entails_{\mathsf{S5}_{I}^{C}} C_Gp,
\]
but no finite subset of the premise set entails \(C_Gp\).  Consequently the
finitary sheet can be weakly complete for valid formulas but not strongly
complete for arbitrary premise sets.  The group-size hypothesis is essential:
for one \(\mathsf{S5}\) agent the iterations collapse.
\end{hbtheorem}

\begin{hbtheorem}{thm:p08-muddy-children}{cited=FaginHalpernMosesVardi1995}{Muddy children under a fixed public protocol}
Suppose \(n\) children see all other foreheads but not their own, exactly
\(k\ge1\) are muddy, the truthful statement that at least one child is muddy is
public, questions are synchronous, and every negative answer is truthful and
public.  The muddy children first know that they are muddy at round \(k\),
after \(k-1\) public rounds of all-negative answers.  The conclusion depends on
all of these information and timing hypotheses; it is not a consequence of the
bare fact that \(k\) children are muddy.
\end{hbtheorem}

\section[Public announcements]{Public announcements and dynamic information}

\begin{hbdefinition}{def:p08-public-announcement}{Truthful public-announcement update}
For an epistemic structure \(\mathcal M\), let
\(\mathcal M\mathbin{\upharpoonright}\chi\) be the substructure induced by the worlds where
\(\chi\) is true.  The public-announcement modality has the clause
\[
 \mathcal M,w\entails[!\chi]A
 \quad\text{ iff }\quad
 \bigl(\mathcal M,w\not\entails\chi\ \lor\
       \mathcal M\mathbin{\upharpoonright}\chi,w\entails A\bigr),
\]
and \(\langle!\chi\rangle A\defeq\neg[!\chi]\neg A\).  The update is a
model restriction, not a change to the valuation at surviving worlds.  Its
standing and vacuous-failure clause distinguish announcement from ordinary
knowledge and are used by the reduction theorem and the unsuccessful-update
example.  General event-model product update is broader and is not identified
with this operation \cite{Plaza1989,BaltagMossSolecki1998}.
\end{hbdefinition}

\begin{hbtheorem}{thm:p08-pal-reduction}{cited=Plaza1989}{Reduction of public announcements}
Over the exact finite-agent \(\mathsf{S5}_{I}\) base fixed by
\hyperref[sheet:hilbert-s5n]{\texttt{sheet:hilbert-s5n}}, every public-announcement
formula is provably equivalent, using
\hyperref[sheet:hilbert-pal]{\texttt{sheet:hilbert-pal}}, to an
announcement-free epistemic formula.  Hence the reduction calculus is sound
and complete for the update semantics of
Definition~\ref{def:p08-public-announcement}.  This does not reduce arbitrary
action-model logic to the same static language
\cite{VanDitmarschVanDerHoekKooi2007}.
\end{hbtheorem}

\begin{systementry}{entry:p08-pal}{Public announcement logic PAL}
  \index{public announcement logic}
  \entrylanguage{The multi-agent epistemic language extended by
    \([!\chi]A\), with \(\langle!\chi\rangle A\) as the executable dual.}
  \entrysemantics{Truth in the induced submodel
    \(\mathcal M\mathbin{\upharpoonright}\chi\) when the truthful announcement
    precondition \(\chi\) holds; a false precondition makes the box formula
    vacuously true.}
  \entrycalculus{sheet:hilbert-pal}
  \entryadequacy{cited=Plaza1989}{The reduction calculus is sound and complete
    for public-announcement validity over that exact \(\mathsf{S5}_{I}\) base.}
  \entrydecidability{cited=VanDitmarschVanDerHoekKooi2007}{Effective reduction
    to the decidable finite-agent \(\mathsf{S5}\) base yields decidability.}
  \entrycharacteristic{Truthful model restriction can change what agents know;
    a true announcement can therefore fail to remain true after its own update.}
  \entrytranslations{Every PAL formula reduces to an announcement-free
    epistemic formula; this does not eliminate arbitrary event-model updates.}
  \entrysources{Plaza1989,BaltagMossSolecki1998,VanDitmarschVanDerHoekKooi2007}
\end{systementry}

\begin{hbexample}{def:p08-unsuccessful-announcement}{An unsuccessful truthful update}
Call \(\chi\) \emph{successful} when \([!\chi]\chi\) is valid.  This is a
metalinguistic property of an announcement formula, needed to distinguish
preservation of announced content from mere executability; it is not another
modality, differs from truth before an update, and is used in this example.
Let two worlds \(w,v\) be indistinguishable to agent \(i\), with \(p\) true
only at \(w\).  At \(w\), \(\chi=p\land\neg K_ip\) is true.  Restricting to
the \(\chi\)-world leaves only \(w\), where \(K_ip\) is now true, so \(\chi\)
is false after its own truthful announcement.  Thus truth before update does
not entail success \cite{VanDitmarschVanDerHoekKooi2007}.
\end{hbexample}

\begin{table}[!htbp]
\centering
\scriptsize
\begin{tabular}{@{}p{.14\textwidth}p{.29\textwidth}p{.20\textwidth}p{.27\textwidth}@{}}
\toprule
Item & Exact semantic role & Rule surface & Evidence boundary \\
\midrule
\(\mathsf{S5}_{I}\) & componentwise equivalence relations; knowledge &
\hyperref[sheet:hilbert-s5n]{\(\mathsf{S5}_n\) sheet} & inherited adequacy remains sketched \\
\(\mathsf{KD45}_{I}\) & componentwise serial, transitive, Euclidean relations;
belief & \hyperref[sheet:hilbert-kd45n]{\(\mathsf{KD45}_n\) sheet} & cited adequacy; no factivity \\
\(E_G\) & finite conjunction of the \(K_i\) operators &
\hyperref[sheet:hilbert-common-knowledge]{common-knowledge sheet} & eliminable abbreviation \\
  \(C_G\) & finite-path closure of the union of agent relations &
  \(C_GA\to E_G(A\land C_GA)\); from the theorem
  \(B\to E_G(A\land B)\) infer the theorem \(B\to C_GA\) & fixed-point and theorem-only induction rule;
  weak completeness only \\
  \(D_G\) & intersection of agent relations; pooled information &
  Definition~\ref{def:p08-distributed-knowledge} & distinct from \(E_G\) and \(C_G\) \\
  \([!\chi]\) & restriction to worlds satisfying \(\chi\) &
  atom/negation/conjunction/\(K_i\)/composition reduction schemas and
  theorem-only replacement on \hyperref[sheet:hilbert-pal]{the PAL sheet} & cited reduction; not full event-model DEL \\
\bottomrule
\end{tabular}
\caption{Epistemic systems and group operators.  The table displays the
common-knowledge fixed-point axiom and theorem-only induction rule in synopsis
and points to its sheet for the complete calculus; the public-announcement
reduction schemas remain sheet-level pointers.}
\label{tab:p08-epistemic-systems}
\end{table}

\begin{hbdefinition}{def:p08-hybrid-k}{Basic hybrid logic \(K({\char64})\)}
Extend modal propositional syntax by a disjoint set of nominals
\(i,j,\ldots\) and satisfaction formulas \(@_iA\). In a Kripke model, each
nominal is true at exactly one world, and
\(M,w\models@_iA\) iff \(M,w_i\models A\), where \(w_i\) is the world named by
\(i\). The name \(K(@)\) is necessary to fix the basic normal modal language
with nominals and satisfaction operators, excluding binders and the global
modality. It stands as a formal language/semantics, differs from ordinary
modal \(K\) by internal world names and jumps, and is used by the labelled
tableau on
\hyperref[sheet:tableaux-hybrid-k]{\texttt{sheet:tableaux-hybrid-k}}
\cite{ArecesTenCate2007}.
\index{hybrid logic}\index{nominal!hybrid logic}\index{satisfaction operator}
\end{hbdefinition}

\begin{hbtheorem}{thm:p08-hybrid-tableau}{cited=ArecesTenCate2007}{Signed-tableau adequacy for basic \(K({\char64})\)}
For the language and semantics of
Definition~\ref{def:p08-hybrid-k}, the complete signed calculus on
\hyperref[sheet:tableaux-hybrid-k]{\texttt{sheet:tableaux-hybrid-k}} has an
open saturated branch from \(@_iA\), for fresh \(i\), exactly when \(A\) is
satisfiable.  Equivalently, every tableau rooted at \(\neg@_iA\) closes exactly
when \(A\) is valid.  The Boolean, modal, satisfaction, and nominal-identity
rules printed on the sheet are all part of the claim.  Binders and the global
modality are outside this exact language and require additional rules.
\end{hbtheorem}

\begin{systementry}{entry:p08-hybrid-k}{Basic hybrid modal logic \(K(@)\)}
  \entrylanguage{Propositional modal \(K\) with nominals and satisfaction
    operators \(@_i\), excluding binders and the global modality.}
  \entrysemantics{Kripke models in which each nominal names exactly one world;
    \(@_iA\) evaluates \(A\) at that named world.}
  \entrycalculus{sheet:tableaux-hybrid-k}
  \entryadequacy{cited=ArecesTenCate2007}{The signed tableau is sound and
    complete for satisfiability and validity, as stated in
    Theorem~\ref{thm:p08-hybrid-tableau}.}
  \entrydecidability{cited=ArecesTenCate2007}{Closure over the finite
    subformula and nominal surface yields a terminating decision procedure for
    the basic language.}
  \entrycharacteristic{Nominals internalize world names and \(@_i\) internalizes
    satisfaction at the named world; neither is definable in ordinary modal
    \(K\) over arbitrary models.}
  \entrytranslations{The standard translation sends a nominal to equality with
    its named-world constant and sends \(@_iA\) to the translation of \(A\)
    evaluated at that constant.}
  \entrysources{Blackburn2000,ArecesTenCate2007}
\end{systementry}

\paragraph{Source and status boundary.}
The exact inherited chapter above supplies the LTL and multi-agent
\(\mathsf{S5}_{I}\) definitions, calculi, adequacy statements and inline
proof sketches, examples, common-knowledge semantics, and all three source
exercises at lines 13340--13582.  CTL, CTL*, modal \(\mu\), tense logic,
\(\mathsf{KD45}_{I}\), the common-knowledge calculus, distributed
knowledge, public announcements, hybrid logic, and the apparatus below are
separate cited or explicitly calculated additions; none inherits a
source-proof length denominator.  The modal countermodel in
Figure~\ref{fig:p08-modal-countermodels}, the LTL timeline in
Figure~\ref{fig:p08-ltl-path}, and the explicit
everyone-without-common-knowledge model in
Example~\ref{def:p08-everyone-not-common} neither replace nor alter the
exact source examples or exercises.
% END HISTORICAL DERIVATIVE 610-774
